\documentclass[10pt,a4paper]{amsart}
\usepackage[margin=19mm]{geometry}
\usepackage{amsmath,amssymb,mathtools}
\usepackage[scr=rsfs,cal=euler]{mathalfa}
\usepackage{xfrac}
\usepackage[unicode,hidelinks,bookmarks=false]{hyperref}
\newcommand{\ie}{i.e.\ }
\newcommand{\cf}{cf.\ }
\newcommand{\resp}{respectively\ }
\newcommand{\Subsection}{\subsection}
\providecommand{\nobreakdash}{\nobreak\hbox{-}}
\setlength{\emergencystretch}{2em}
\multlinegap0pt
\usepackage{amssymb}
\usepackage{mathtools}
\usepackage{xcolor}
\usepackage{algorithm}\usepackage{algpseudocode}
\newtheorem{corollary}{Corollary}
\newtheorem{lemma}{Lemma}
\newtheorem{theorem}{Theorem}
\usepackage{cleveref}
\crefname{corollary}{Corollary}{Corollarys}
\crefname{lemma}{Lemma}{Lemmas}
\crefname{theorem}{Theorem}{Theorems}
\newcommand{\GL}{\mathcal{G}_L}
\newcommand{\R}{\mathcal{R}}
\newcommand{\stack}{\Pi}
\title{The Extremum Stack is a Minimal Sufficient Statistic for Rate-Independent Functionals:\\ A Kolmogorov Complexity Characterisation}
\author{Piotr Frydrych}
\date{Source: arXiv:2605.18885v1 (CC BY 4.0)}
\begin{document}
\maketitle

\noindent\emph{Source and licence.} The Extremum Stack is a Minimal Sufficient Statistic for Rate-Independent Functionals:\\ A Kolmogorov Complexity Characterisation. Version arXiv:2605.18885v1 (CC BY 4.0), distributed by arXiv under the Creative Commons Attribution 4.0 International licence, http://creativecommons.org/licenses/by/4.0/, as recorded in the arXiv licence metadata for this exact version and shown on the abstract page https://arxiv.org/abs/2605.18885v1. Full licence notice: \texttt{licenses/arxiv-2605.18885-CC-BY-4.0.txt}.\par\smallskip

\noindent\textit{Editorial scope.} Counted unit: the article's theorem thm:main (source0.tex lines 412--424). The article's complete original proof of it is delivered with it (source0.tex lines 426--444, one \texttt{proof} environment of 19 lines). No mathematics is written, repaired or re-derived: the statement and the proof are the article's own lines, in the article's own notation, and no retained line is substituted or re-flowed.  Retained uncounted as the source-local context the proof needs: lines 313--316; lines 384--392.  Nothing was omitted from any retained span, so the package carries no omission record.  No retained line contains a citation, so the wrapper carries no bibliography and no bibliography entry.  No retained line contains a figure environment, an image-inclusion command or drawing code, so no image is delivered and the package declares no asset dependency.  Editorial formatting only, in addition: an AMS \texttt{amsart} preamble that restates the article's own declarations and macros used by the retained lines, namely the environments \texttt{corollary}, \texttt{lemma}, \texttt{theorem} on separate counters, together with the standard packages those lines need.  The retained environments print in this document as Corollary 1, Lemma 1, Theorem 1 in order of appearance, whereas the article prints the same items with the numbering of its own full document.  The complete native source of the article as distributed in the arXiv e-print for this exact version is bundled unchanged as \texttt{sources/200880/source0.tex}.
\begin{corollary}[Upper bound]
\label{cor:upper}
$K_{\R}(u_{0:n}) \leq K(\stack_n) + O(1)$.
\end{corollary}

\begin{lemma}[Stack reconstruction]
\label{lem:recon}
Let $\mathcal{S}$ be any set and
$R: \GL^{n+1} \to \mathcal{S}$ a representation such that
every $F \in \R$ is computable from $R(u_{0:n})$.
Then there exists a computable function
$\Phi: \mathcal{S} \to (\GL \times \GL)^*$
such that $\Phi(R(u_{0:n})) = \stack_n$.
\end{lemma}

\begin{theorem}[Minimality of the extremum stack]
\label{thm:main}
Let $u_{0:n} \in \GL^{n+1}$ and $\stack_n$ its extremum stack.
Then:
\begin{equation}
  K(\stack_n) - O(1)
  \;\leq\; K_{\R}(u_{0:n})
  \;\leq\; K(\stack_n) + O(1).
  \label{eq:main}
\end{equation}
The extremum stack is, up to an additive constant, the
shortest program that answers every $\R$-query.
\end{theorem}

\begin{proof}
\textbf{Lower bound.}
By \cref{lem:recon}, there exists a computable $\Phi$ with
$\Phi(R^*(u_{0:n})) = \stack_n$, where $R^*$ is the
$K_{\R}$-optimal representation.
Therefore:
\[
  K(\stack_n) \leq K(R^*(u_{0:n})) + O(1) = K_{\R}(u_{0:n}) + O(1),
\]
since $\stack_n$ is computable from the optimal representation
by a program of fixed length $O(1)$ (the description of $\Phi$).
Rearranging gives the lower bound.

\textbf{Upper bound.}
The upper bound follows directly from \cref{cor:upper}:
since every $F[u](n)$ is computable from $\stack_n$ by a fixed
$O(1)$-length program, the prefix-free invariance theorem gives
$K_{\R}(u_{0:n}) \leq K(\stack_n) + O(1)$.
\end{proof}

\end{document}
